\documentclass[12pt]{article}
\usepackage{ctex}
\input{../../style/math_commands.tex}
\input{../../style/note.sty}
\usepackage[colorlinks=true,linkcolor=blue,urlcolor=blue]{hyperref}

\title{Lecture 1}

\begin{document}
\begin{center}
  {\LARGE \bf Lecture 1 计算机核心数学根基}\\
  {\large 周建宇}
\end{center}


\section{数学建模(DAY 1)}

\[
  \text{问题}\rightarrow \text{模型}\rightarrow \text{算法}\rightarrow \text{代码}
\]
\subsection{KMP 算法}
在字符串匹配问题中，
主串（文本）记为 $T$，其长度为 $n$；
模式串（待查找的子串）记为 $P$，其长度为 $m$。
字符串 $T$ 的第 $i$ 个字符记为 $T[i]$，模式串 $P$ 的第 $j$ 个字符记为 $P[j]$。
区间 $T[i..j]$ 表示从 $T[i]$ 到 $T[j]$ 的连续子串。

问题定义：给定主串 $T[1..n]$ 和模式串 $P[1..m]$，求所有起始位置 $i$，使得
\[
  T[i..i+m-1] = P[1..m],\quad 1 \le i \le n-m+1.
\]

这里的“匹配”表示模式串 $P$ 的每个字符与对应主串子串的字符一一相等；“失配”表示当前比较的字符不相等。

这是一个典型的字符串匹配问题：从问题出发，我们先建立模式串的前缀函数模型，再用该模型指导匹配过程。

\begin{defn}
  KMP 算法（Knuth--Morris--Pratt 算法）是一种字符串模式匹配算法。它通过预处理模式串得到前缀函数 $\pi$，记录每个位置最长相等真前缀和真后缀长度，从而在失配时无需回退主串指针，实现 $O(n+m)$ 时间匹配。
\end{defn}
\begin{defn}[真前缀与真后缀]
  给定字符串 $P[1..q]$：
  \begin{itemize}
    \item \newterm{真前缀（proper prefix）}：$P$ 的前缀但不包含整个字符串本身，即形如 $P[1..k]$ 且 $0 \le k < q$ 的子串；
    \item \newterm{真后缀（proper suffix）}：$P$ 的后缀但不包含整个字符串本身，即形如 $P[q-k+1..q]$ 且 $0 \le k < q$ 的子串。
  \end{itemize}
\end{defn}

\begin{defn}[前缀函数]
  模式串 $P[1..m]$ 的前缀函数 $\pi[1..m]$ 定义为：每个位置 $q$ 处的值等于 $P[1..q]$ 的最长相等真前缀与真后缀的长度，即
  \[
    \pi[1] = 0,
  \]
  \[
    \pi[q] = \max\{k \mid 0 \le k < q,\ P[1..k] = P[q-k+1..q]\},\quad 2 \le q \le m.
  \]
\end{defn}
例如，令模式串
\[
  P = \text{ababaca},\quad m = 7.
\]
则各位置的前缀函数为：
\[
  \pi = [0, 0, 1, 2, 3, 0, 1],
\]
其中：
\begin{itemize}
  \item $\pi[3]=1$，因为 $P[1..1]=\text{a}$ 与 $P[3..3]=\text{a}$，且没有更长的真前缀与真后缀匹配；
  \item $\pi[5]=3$，因为 $P[1..3]=\text{aba}$ 与 $P[3..5]=\text{aba}$；
  \item $\pi[6]=0$，因为 $P[1] \neq P[6]$，此时没有非空真前缀与真后缀相同；
  \item $\pi[7]=1$，因为 $P[1..1]=\text{a}$ 与 $P[7..7]=\text{a}$。
\end{itemize}
匹配过程：令主串 $T[1..n]$，模式串 $P[1..m]$，当前已匹配长度 $q=0$。
遍历 $i=1\dots n$ 表示按顺序扫描主串 $T$ 的第 $i$ 个字符 $T[i]$，在每一步将它与模式串中下一个待比较字符 $P[q+1]$ 进行比较。
若失配，则
\[
  q \leftarrow \pi[q-1]
\]
直到重新匹配或退到 $0$。

例如，令主串和模式串为：
\[
  T = \text{abababaca},\qquad P = \text{ababaca}.
\]
已知 $\pi = [0,0,1,2,3,0,1]$。
前 6 步匹配过程如下：
\begin{itemize}
  \item $i=1$：$T[1]=\text{a}$ 与 $P[1]=\text{a}$ 匹配，$q=1$；
  \item $i=2$：$T[2]=\text{b}$ 与 $P[2]=\text{b}$ 匹配，$q=2$；
  \item $i=3$：$T[3]=\text{a}$ 与 $P[3]=\text{a}$ 匹配，$q=3$；
  \item $i=4$：$T[4]=\text{b}$ 与 $P[4]=\text{b}$ 匹配，$q=4$；
  \item $i=5$：$T[5]=\text{a}$ 与 $P[5]=\text{a}$ 匹配，$q=5$；
  \item $i=6$：$T[6]=\text{b}$ 与 $P[6]=\text{c}$ 失配，于是 $q \leftarrow \pi[5] = 3$，接着比较 $P[4]=\text{b}$ 与 $T[6]=\text{b}$，继续匹配，得到 $q=4$；
  \item $i=7$：$T[7]=\text{a}$ 与 $P[5]=\text{a}$ 匹配，$q=5$；
  \item $i=8$：$T[8]=\text{c}$ 与 $P[6]=\text{c}$ 匹配，$q=6$；
  \item $i=9$：$T[9]=\text{a}$ 与 $P[7]=\text{a}$ 匹配，$q=7=m$，找到一个完整匹配，起始位置为 $i-m+1 = 3$。
\end{itemize}
这样，KMP 利用前缀函数避免了回退主串指针，直接从模式串已知的最长可重用前缀长度继续比较。

计算前缀函数和查找匹配的伪代码如下：
\begin{verbatim}
function compute_prefix(P):
    m = len(P)
    pi = array[1..m]
    pi[1] = 0
    k = 0
    for q = 2 to m:
        while k > 0 and P[k+1] != P[q]:
            k = pi[k]
        if P[k+1] == P[q]:
            k = k + 1
        pi[q] = k
    return pi

function kmp_search(T, P):
    n = len(T)
    m = len(P)
    pi = compute_prefix(P)
    q = 0
    results = []
    for i = 1 to n:
        while q > 0 and P[q+1] != T[i]:
            q = pi[q]
        if P[q+1] == T[i]:
            q = q + 1
        if q == m:
            results.append(i - m)
            q = pi[q]
    return results
\end{verbatim}

\begin{note}
  与 KMP 相关的字符串进阶数据结构与算法：
  \begin{itemize}
    \item \newterm{Aho-Corasick 自动机}：多模式串匹配算法。把全部模式串建成 Trie 树，再仿照 KMP 的前缀函数为每个节点建立失配指针（fail 指针），扫描主串一次即可同时匹配所有模式串，总复杂度 $O(n + m + \text{匹配数})$（$m$ 为所有模式串长度之和），常用于敏感词过滤、多关键字检索。
    \item \newterm{后缀数组（Suffix array, SA）}：把字符串的所有后缀按字典序排序得到的数组，配合最长公共前缀数组（LCP 数组）可高效解决子串查找、最长重复子串、不同子串计数等问题，倍增法构建复杂度 $O(n\log n)$。
    \item \newterm{后缀自动机（Suffix automaton, SAM）}：能识别一个字符串的所有子串的最小确定性有限自动机，状态数与转移数均为 $O(n)$，可以线性时间在线构造，用于统计子串出现次数、求最长公共子串等。
    \item \newterm{后缀树（Suffix tree）}：把所有后缀压缩存储的字典树，$O(n)$ 个节点，可在 $O(|P|)$ 时间内判断模式串 $P$ 是否为主串的子串。
    \item \newterm{滚动哈希（Rolling hash）}：把字符串视作一个进制数并取模，用 $O(1)$ 时间从上一窗口的哈希值推出相邻窗口的哈希值，从而快速比较子串是否相等；理论上有碰撞风险，实践中可用双哈希降低碰撞概率。
    \item \newterm{倒排索引（Inverted index）}：从词项映射到包含该词的文档列表的索引结构（词典 $\to$ 倒排表），与"文档 $\to$ 词项"的正排索引相对，是搜索引擎全文检索的核心数据结构。
    \item \newterm{Shingling}：将文档切分为连续的 $k$ 个 token 组成的集合（$k$-gram），用集合的 Jaccard 相似度度量两份文档的相似程度。
    \item \newterm{MinHash}：对集合元素施加多种随机哈希，取每组哈希值的最小值作为文档签名，使得签名相等的概率恰好等于集合的 Jaccard 相似度，从而在大规模语料上近似估计文档相似度（网页去重）。
    \item \newterm{SimHash}：把文本中每个特征词哈希为二进制向量，按权重累加后取符号得到定长指纹，两个指纹的汉明距离小则认为文档相似；适合海量网页去重。
  \end{itemize}
\end{note}

\begin{defn}[Dijkstra 算法]
  求解非负权图单源最短路径的经典算法。维护集合 $S$（已确定最短距离的顶点）和距离数组 $d$，每轮从未确定顶点中选出 $d$ 最小的顶点 $u$ 加入 $S$，并对其所有出边做松弛操作；用二叉堆优化后复杂度为 $O((n+m)\log n)$。
\end{defn}

\begin{defn}[松弛操作（relaxation）]
  对边 $(u,v)$，若
  \[
    d[v] > d[u] + w(u,v),
  \]
  则更新 $d[v] \leftarrow d[u] + w(u,v)$。松弛是 Dijkstra、Bellman-Ford、SPFA 等最短路径算法共用的核心步骤：用一条"可能更短的路"去改善当前已知的最短距离估计。
\end{defn}

Dijkstra 算法的伪代码如下（$\mathrm{PQ}$ 为小根堆，键为距离）：
\begin{verbatim}
function dijkstra(s):
    for v in V: d[v] = INF
    d[s] = 0
    PQ = min-heap containing (0, s)
    while PQ not empty:
        (du, u) = PQ.extract_min()
        if du != d[u]: continue      // 过期条目
        for each edge (u, v, w):
            if d[v] > d[u] + w:      // 松弛操作
                d[v] = d[u] + w
                PQ.insert((d[v], v))
    return d
\end{verbatim}
\[
  \text{问题}\begin{cases}
    \text{生活问题} \\
    \text{科学问题} \\
    \text{可建模的科学问题}    \begin{cases}
                         \text{假设} \\
                         \text{规模}
                       \end{cases}
  \end{cases}
\]
\begin{note}
  作业：报告，可AI生成，回顾本科所学的知识，有哪些知识，经过这五天的学习，哪些知识有了更深入的理解，前因后果，核心理解。
\end{note}
\begin{exam}
  计算机如何计算$\binom{m}{n}=\frac{m!}{n!(m-n)!}?=\frac{\prod_{i=1}^{n} (m-i+1)}{\prod_{i=1}^{n} i}=\prod_{i=1}^{n} \frac{m-i+1}{i}$， $n$个连续自然数一定会被$n$整除,or利用组合数公式$\binom{m}{n}=\binom{m-1}{n-1}+\binom{m-1}{n}$。

  \newterm{思路}：直接计算 $m!/(n!(m-n)!)$ 在 $m$ 较大时很快溢出；改按 $\prod_{i=1}^{n}\frac{m-i+1}{i}$ 逐步相乘可避免中间值爆炸（每一步结果都是整数，因为 $n$ 个连续自然数必含 $n$ 的因子）。也可以按帕斯卡递推 $\binom{m}{n}=\binom{m-1}{n-1}+\binom{m-1}{n}$ 用 $O(nm)$ 的动态规划填表，且中间值保持在 $\binom{m}{n}$ 同量级内，适合精确大数计算。

\end{exam}
\subsection{等价变化}
\begin{exam}
  $m$个完全的小球放进$n$个完全相同的盘子里?
  \[
    F_{m,n}=\begin{cases}
      1,                                                       & n=1 \text{或} m=0 \\
      F_{m,m}                                                  & n>m              \\
      F_{m,n-1}(\text{有空盘})+F_{m-n,n}(\text{无空盘，每个盘子至少有一个小球}), & n>1 \text{且} m>0
    \end{cases}
  \]
  \newterm{思路}：由于球和盘都不可区分，只需关心"每盘装几个"的非增整数分拆（partition）。递推的关键是枚举是否允许空盘：若至少有一个空盘，则去掉一个盘子化为 $F_{m,n-1}$；若没有空盘，则先从每个盘子拿掉一个小球，化为 $F_{m-n,n}$。边界 $n>m$ 时空盘不可避免，等价于 $F_{m,m}$。这是"整数分拆数"模型的等价变化。
\end{exam}
\begin{exam}[错位排序]
  $n$个盒子有不同标签$1,2,\ldots,n$，$n$个不同的小球$1,2,\ldots,n$，有不同标签，放进盒子里，每个盒子有一个小球，问有多少种放法？$n!$种放法。
  若要求盒子的标签与小球的标签不同，问有多少种放法？设$f(n)$为所求的放法数，则有递推关系：
  \[
    f(n)=(n-1)(f(n-1)+f(n-2)),\quad n\ge 3
  \]
  \newterm{思路（容斥 / 第一球落位分析）}：错位排列的经典证明：考虑小球 $n$，它可放在 $n-1$ 个非盒子 $n$ 的位置之一。若小球 $n$ 放在盒子 $k$ 且小球 $k$ 放在盒子 $n$，则剩下 $n-2$ 个球全错位，共 $f(n-2)$ 种；若小球 $k$ 不放盒子 $n$，则把"盒子 $n$ 必须避开小球 $k$"纳入错位要求，等价于 $n-1$ 个元素的错位排列 $f(n-1)$ 种。于是 $f(n)=(n-1)(f(n-1)+f(n-2))$。也可用容斥原理 $f(n)=n!\sum_{k=0}^{n}\frac{(-1)^k}{k!}$。
\end{exam}
\begin{exam}
  $n$个节点的二叉树有多少种？设$f(n)$为所求的二叉树数，则有递推关系：
  \[
    f(n)=\sum_{i=0}^{n-1} f(i)f(n-i-1),\quad n\ge 1
  \]
  \newterm{思路}：按根节点划分左右子树。固定根后，左子树含 $i$ 个节点、右子树含 $n-i-1$ 个节点，左右独立，故对 $i$ 求和得卷积递推 $f(n)=\sum_{i=0}^{n-1}f(i)f(n-i-1)$，即卡特兰数的递推定义，见下。
\end{exam}
\begin{note}
  该递推给出的是第 $n$ 个卡特兰数(Catalan number) $C_n$（记 $f(n)=C_n$）。
  闭式为
  \[
    C_n=\frac{1}{n+1}\binom{2n}{n},\quad C_0=1.
  \]
  生成函数为
  \[
    C(x)=\sum_{n\ge0}C_nx^n=\frac{1-\sqrt{1-4x}}{2x}.
  \]
  渐近行为为 $C_n\sim\dfrac{4^n}{\sqrt{\pi}\,n^{3/2}}$。卡特兰数计数多种组合结构，例如二叉树、合法括号序列、非交错路径等。
\end{note}
\begin{defn}[Floyd-Warshall 算法]
  求解全源最短路径的动态规划算法。设 $d^{(k)}_{ij}$ 为只允许经过 $\{1,\dots,k\}$ 作为中间点时 $i$ 到 $j$ 的最短路，递推式
  \[
    d^{(k)}_{ij} = \min\{d^{(k-1)}_{ij},\ d^{(k-1)}_{ik}+d^{(k-1)}_{kj}\},
  \]
  复杂度 $O(n^3)$，允许负权边但不允许负环，常用于传递闭包等图论问题。伪代码如下：
  \begin{verbatim}
  function floyd_warshall(W):          // W 为边权矩阵
      d = W
      for k = 1 to n:
          for i = 1 to n:
              for j = 1 to n:
                  d[i][j] = min(d[i][j], d[i][k] + d[k][j])
      return d
  \end{verbatim}
\end{defn}

\begin{defn}[Bellman-Ford 算法]
  求解单源最短路径，可处理负权边。对全部边执行 $n-1$ 轮松弛，第 $k$ 轮结束后得到的是经过至多 $k$ 条边的最短路；若第 $n$ 轮仍能松弛，则存在负环。复杂度 $O(nm)$。伪代码如下：
  \begin{verbatim}
  function bellman_ford(s):
      for v in V: d[v] = INF
      d[s] = 0
      for k = 1 to n-1:
          for each edge (u, v, w):
              if d[v] > d[u] + w: d[v] = d[u] + w   // 松弛
      for each edge (u, v, w):                        // 检测负环
          if d[v] > d[u] + w: return "存在负环"
      return d
  \end{verbatim}
\end{defn}

\begin{defn}[SPFA]
  Bellman-Ford 的队列优化版本，只对"距离被更新过的顶点"松弛其出边，通常很快但最坏仍为 $O(nm)$，也可用于检测负环。
\end{defn}

\begin{defn}[维特比算法（Viterbi algorithm）]
  在隐马尔可夫模型（HMM）中求解最可能隐状态序列的动态规划算法，即给定观测序列求使后验概率最大的状态路径；复杂度 $O(T N^2)$（$T$ 为序列长度，$N$ 为状态数），广泛用于词性标注、语音识别、纠错解码等。
\end{defn}

\begin{defn}[逆序对的排列]
  若排列 $\pi=(\pi_1,\dots,\pi_n)$ 中存在 $i<j$ 且 $\pi_i>\pi_j$，则 $(\pi_i,\pi_j)$ 是一个逆序对；逆序对数刻画排列的"乱序程度"，并与冒泡排序的交换次数、排列的奇偶性密切相关。
\end{defn}
\subsection{重要假设}
\begin{exam}
  \newterm{贝特朗悖论（Bertrand paradox）}：单位圆上任取一条弦，问弦长 $d \ge \sqrt{3}$（即不小于内接正三角形边长）的概率是多少？因为"随机取弦"的等可能性定义方式不同，会得到不同答案：
  \begin{itemize}
    \item $\frac{1}{3}$：随机取两个端点（端点在圆周上均匀分布）——对应圆心角之差在 $[120^\circ,240^\circ]$ 或等价的角范围；
    \item $\frac{1}{2}$：随机取弦的中点沿某条直径均匀分布——中点必须落在距圆心 $\le \frac12$ 的范围内；
    \item $\frac{1}{4}$：随机取弦的中点在整个圆内均匀分布——中点必须落在半径为 $\frac12$ 的同心圆内。
  \end{itemize}
  该悖论说明：同一几何问题在"均匀"的不同理解下概率可以不同，数学建模必须先明确假设。
\end{exam}
\begin{defn}[K-means 算法]
  给定 $n$ 个样本和簇数 $k$，把样本划分到 $k$ 个簇中，使簇内平方和
  \[
    \sum_{i=1}^{k}\sum_{x\in C_i}\lVert x-\mu_i\rVert^2
  \]
  最小（即最小化类内距离，等价于间接最大化类间距离）。交替进行"分配：每个点归到最近的质心"与"更新：重新计算质心"两步直至收敛；是 NP-hard 问题的启发式求解，复杂度 $O(nkt)$，对初始质心和 $k$ 的选择敏感。伪代码如下：
  \begin{verbatim}
  function kmeans(X, k):
      mu = random_k_centroids(X)
      repeat until centroids stop changing:
          C = empty assignment for each x
          for each x in X:
              assign x to cluster of nearest centroid mu
          for i = 1 to k:
              mu[i] = mean of all points in cluster i
      return clusters
  \end{verbatim}
\end{defn}

\begin{defn}[距离度量]
  距离度量的选择影响聚类结果：常用欧氏距离 $\lVert x-y\rVert_2$、曼哈顿距离 $\lVert x-y\rVert_1$、余弦相似度 $\frac{x\cdot y}{\lVert x\rVert\lVert y\rVert}$ 等，距离定义本身就是建模中的关键假设。
\end{defn}

\begin{defn}[层次聚类]
  不预设簇数，自底向上（凝聚式）每次合并距离最近的两个簇，或自顶向下（分裂式）每次拆分，最终得到一棵聚类树（树状图），可在任意高度切割得到不同粒度的簇。凝聚式层次聚类的伪代码如下：
  \begin{verbatim}
  function agglomerative_clustering(X):
      each point is its own cluster
      while more than one cluster remains:
          merge the two clusters closest under linkage
          record merge in dendrogram
      return dendrogram
  \end{verbatim}
\end{defn}
\begin{defn}[DBSCAN]
  基于密度的聚类算法，不需预设簇数。给定邻域半径 $\varepsilon$ 和最小点数 $\mathrm{MinPts}$：若某点 $\varepsilon$-邻域内点数 $\ge \mathrm{MinPts}$ 则为核心点；核心点及其密度可达的点构成一个簇，其余点记为噪声。能发现任意形状的簇并识别离群点，配合空间索引复杂度约 $O(n\log n)$。伪代码如下：
  \begin{verbatim}
  function dbscan(X, eps, MinPts):
      label each point as unvisited
      for each unvisited point p:
          N = points within eps of p
          if |N| < MinPts: mark p as noise
          else:
              start new cluster C with p
              expand C to all density-reachable points
      return clusters
  \end{verbatim}
\end{defn}
\section{计算机硬件与问题规模(DAY 2)}
\subsection{问题规模}
\begin{exam}
  算法：基于比较的排序。给定 $n$ 个元素，比较类排序算法的最坏时间复杂度下界为 $\Omega(n\log n)$（基于决策树）。常见排序算法及复杂度：
  \begin{itemize}
    \item \newterm{冒泡排序}：反复扫描序列，相邻逆序元素就交换，每趟把当前最大元素"冒泡"到末尾；$O(n^2)$，最好情况 $O(n)$。
    \item \newterm{插入排序}：将每个元素插入到已排序前缀的正确位置；$O(n^2)$，对近似有序序列接近 $O(n)$，稳定。
    \item \newterm{快速排序}：分治 + 划分，选枢轴把序列分成左右两部分递归排序；平均 $O(n\log n)$，最坏 $O(n^2)$（见下），原地排序但递归栈占 $O(\log n)$ 空间，不稳定。
    \item \newterm{归并排序}：分治，先递归排序左右两半再线性归并；稳定，$O(n\log n)$，但需要 $O(n)$ 额外空间。归并排序的合并阶段伪代码如下：
  \end{itemize}
  \begin{verbatim}
  function merge_sort(A, l, r):
      if l >= r: return
      mid = (l + r) / 2
      merge_sort(A, l, mid); merge_sort(A, mid+1, r)
      merge two sorted halves in O(r-l+1)
  \end{verbatim}
\end{exam}
\begin{thm}[比较排序下界]
  任何基于比较的排序算法的最坏时间复杂度至少为
  \[
    \Omega(n\log n),
  \]
  精确地，$T(n) \ge \lceil\log_2(n!)\rceil$。
\end{thm}
\begin{proof}
  每次比较把可能的排列按结果分为两支，形成一棵决策树；排序算法必须能区分全部 $n!$ 种排列，故决策树的叶子数 $\ge n!$，高度至少 $\lceil\log_2(n!)\rceil$。由 Stirling 公式 $n!\sim\sqrt{2\pi n}(n/e)^n$ 得 $\log_2(n!)=\Theta(n\log n)$。
\end{proof}
\begin{thm}[主定理（Master Theorem）]
  设递归式
  \[
    T(n)=aT(n/b)+f(n),
  \]
  其中 $a\ge 1$，$b>1$ 为常数，且 $f(n)$ 为渐近正函数。则：
  \begin{itemize}
    \item 若 $f(n)=O(n^{\log_b a-\varepsilon})$，其中 $\varepsilon>0$ 为任意小的常数，则 $T(n)=\Theta(n^{\log_b a})$；
    \item 若 $f(n)=\Theta(n^{\log_b a}\log^k n)$，其中 $k\ge 0$，则 $T(n)=\Theta(n^{\log_b a}\log^{k+1} n)$；
    \item 若 $f(n)=\Omega(n^{\log_b a+\varepsilon})$，其中 $\varepsilon>0$，且满足正则条件 $a f(n/b)\le c f(n)$（对足够大的 $n$），则 $T(n)=\Theta(f(n))$。
  \end{itemize}
\end{thm}
\begin{proof}[主定理证明思路（递归树）]
  把递归式 $T(n)=aT(n/b)+f(n)$ 展开为递归树：树高 $h=\log_b n$，第 $j$ 层的代价为 $a^j f(n/b^j)$，叶节点数为 $a^h = n^{\log_b a}$，每个叶节点代价 $O(1)$。总代价
  \[
    T(n)=\sum_{j=0}^{\log_b n-1} a^j f(n/b^j) + \Theta(n^{\log_b a}).
  \]
  \begin{itemize}
    \item 情形 1：$f(n)=O(n^{\log_b a-\varepsilon})$ 时，各层代价按几何级数递减，和由叶节点主导，故 $T(n)=\Theta(n^{\log_b a})$；
    \item 情形 2：$f(n)=\Theta(n^{\log_b a}\log^k n)$ 时，每层代价近似相等，共 $\log_b n$ 层，故 $T(n)=\Theta(n^{\log_b a}\log^{k+1} n)$；
    \item 情形 3：$f(n)=\Omega(n^{\log_b a+\varepsilon})$ 且正则条件保证非叶层代价递增，和由根主导，故 $T(n)=\Theta(f(n))$。
  \end{itemize}
\end{proof}
\begin{note}
  给出一个“不开辟额外空间”的快速排序实现思路：
  \begin{verbatim}
  def quick_sort(arr, left, right):
      if left >= right:
          return
      pivot = arr[right]
      i = left - 1
      for j in range(left, right):
          if arr[j] <= pivot:
              i += 1
              arr[i], arr[j] = arr[j], arr[i]
      arr[i + 1], arr[right] = arr[right], arr[i + 1]
      mid = i + 1
      quick_sort(arr, left, mid - 1)
      quick_sort(arr, mid + 1, right)
  \end{verbatim}
\end{note}
\begin{note}
  快速排序的平均时间复杂度（均匀性假设下）为 $O(n \log n)$，最坏情况为 $O(n^2)$。若每次划分都把序列分成 $1$ 和 $n-1$ 两部分，就会退化到 $O(n^2)$。例如，若每次都选到最小值或最大值作为枢轴。对于 $n$ 个互异元素，若每一层都只有“选最大”或“选最小”这两种选择，那么总共有 $2^{n-1}$ 种最坏情况。
\end{note}
\begin{defn}[群（group）]
  集合 $G$ 连同二元运算 $\circ$，满足\textbf{封闭性}（$a\circ b\in G$）、\textbf{结合律}（$(a\circ b)\circ c=a\circ(b\circ c)$）、存在\textbf{单位元} $e$（$\forall a\in G,\ a\circ e=e\circ a=a$）、每个元素有\textbf{逆元}（$\forall a\in G,\ \exists a^{-1},\ a\circ a^{-1}=a^{-1}\circ a=e$）。若还满足交换律则称为\textbf{阿贝尔群}。例：整数加法群 $(\mathbb{Z},+)$。
\end{defn}

\begin{defn}[环（ring）]
  集合 $R$ 上定义加法和乘法：$(R,+)$ 构成阿贝尔群，乘法满足结合律，且乘法对加法满足分配律。例：整数环 $\mathbb{Z}$。
\end{defn}

\begin{defn}[域（field）]
  交换环中若每个非零元都有乘法逆元（即 $(R\setminus\{0\},\times)$ 构成阿贝尔群），则称为域。例：$\mathbb{Q},\mathbb{R},\mathbb{C}$ 及有限域 $\mathbb{F}_p$。
\end{defn}

\begin{defn}[格（lattice）]
  偏序集中任意两个元素都有最小上界（并）和最大下界（交）的结构；也可指数论/密码学中的 $\mathbb{R}^n$ 的离散加法子群（如 $\mathbb{Z}^n$），是格基密码的基础。
\end{defn}
\begin{defn}[确定性算法]
  给定相同输入，执行的每一步都完全确定，输出必然相同。
\end{defn}

\begin{defn}[随机算法（randomized algorithm）]
  执行过程中利用随机性作决策的算法。两类典型：
  \begin{itemize}
    \item \newterm{拉斯维加斯算法（Las Vegas algorithm）}：输出一定正确，但运行时间是一个随机变量（可能一直跑下去但概率为 0），如随机化快速排序；
    \item \newterm{蒙特卡洛算法（Monte Carlo algorithm）}：运行时间确定，但输出可能以一定概率出错，可通过重复运行降低错误概率，如 MinHash、Miller-Rabin 素性测试。
  \end{itemize}
  例如对快速排序先做随机打乱（shuffle），即可避免固定有序输入导致的 $O(n^2)$ 最坏退化。
\end{defn}
\begin{note}
  怎样改进快速排序算法，产生$n!$种确定性算法？shuffle，随机打乱输入序列。
\end{note}
\begin{note}
  如何查找 $n$ 个元素中第 $k$ 小的元素？

  一种思路是使用堆排序：先构建一个最大堆，然后依次取出堆顶元素。建堆的时间复杂度为 $O(n)$，每次取出堆顶并调整堆的时间复杂度为 $O(\log n)$，因此取出 $k$ 次的总复杂度为 $O(n+k\log n)$。

  另一种更高效的思路是快速选择算法。其核心思想与快速排序类似：通过一次划分把数组分成两部分。若枢轴位置恰好是 $k$，则直接返回；若 $k$ 小于枢轴位置，则只递归处理左边；若 $k$ 大于枢轴位置，则只递归处理右边。其伪代码如下：
  \begin{verbatim}
  function quick_select(arr, left, right, k):
    p = partition(arr, left, right)
    rank = p - left + 1

    if rank == k:
        return arr[p]
    else if rank > k:
        return quick_select(arr, left, p - 1, k)
    else:
        return quick_select(arr, p + 1, right, k - rank)
  \end{verbatim}
  $T(n)=T(n/2)+O(n)$，因此平均复杂度为 $O(n)$。
\end{note}
\begin{exam}
  如何求逆序对个数？利用归并排序代码。

  对于一个序列 $a_1,a_2,\dots,a_n$，若 $i<j$ 且 $a_i>a_j$，则 $(i,j)$ 就是一个逆序对。可以在归并排序的合并阶段顺带统计逆序对个数：合并时，如果左指针元素大于右指针元素，则右指针元素与左半部分剩余的所有元素构成逆序对。

\end{exam}
\subsection{计算机硬件}
\begin{note}
  为什么实际过程中不常用归并排序而是用快速排序？归并排序需要额外的空间，而快速排序是原地排序。归并排序的时间复杂度为 $O(n \log n)$，空间复杂度为 $O(n)$；快速排序的平均时间复杂度为 $O(n \log n)$，空间复杂度为 $O(\log n)$（递归栈空间）。
\end{note}
\begin{note}
  常见存储单位之间的换算关系如下：8 bit = 1 byte，1 KB = 1024 byte，1 MB = 1024 KB，1 GB = 1024 MB，1 TB = 1024 GB，1 PB = 1024 TB，1 EB = 1024 PB。

  个人电脑内存 4-8 GB，服务器内存 16-64 GB，云服务器内存 128 GB 以上。
\end{note}
\begin{defn}[外部排序（external sort）]
  当数据量超过内存容量时，无法一次全部读入内存排序，于是分块读入、块内排序生成有序"归并段"（run），再通过多路归并合成完整有序序列，是归并排序在磁盘上的推广。

  当单机内存/磁盘都不够用时，需要把计算分发到多台机器，即\newterm{分布式程序}。典型框架：
  \[
    \text{Hadoop}\begin{cases}
      \text{MapReduce} \\
      \text{Spark}
    \end{cases}
  \]
  \newterm{MapReduce}：把计算抽象为 Map（将输入映射为中间键值对）和 Reduce（按键聚合处理）两个阶段，由框架负责调度与容错，适合离线批处理。

  \newterm{Spark}：基于内存的分布式计算引擎，利用 RDD（弹性分布式数据集）在内存中缓存中间结果，避免反复读写磁盘，因此对迭代式计算（如机器学习）远快于磁盘型的 MapReduce。
\end{defn}
\begin{note}
  \[
    \text{主机}\begin{cases}
      \text{CPU} \\
      \text{主存}
    \end{cases}\leftrightarrow_{\text{PCIe}} \text{显卡}\begin{cases}
      \text{GPU} \\
      \text{显存}
    \end{cases}
  \]
\end{note}
\begin{defn}[译码器（decoder）]
  组合逻辑电路，$n$ 位输入产生 $2^n$ 条输出线且恰有一条为有效，用于把地址翻译成对应的存储单元选通信号；HBM 等内存内部依赖译码器完成行列寻址。
\end{defn}

\begin{defn}[HBM（High Bandwidth Memory，高带宽内存）]
  通过 TSV（硅通孔）垂直堆叠多层 DRAM 实现的高带宽内存，显著高于普通 DDR 内存的带宽且功耗更低，常用于 GPU 和 AI 加速卡（如 HBM2/HBM3）。
\end{defn}
\begin{note}
  \[
    \text{CPU}\begin{cases}
      \text{intel(INTC)}(\text{指令，优化})\begin{cases}
                                        \text{家用}\text{u9,u7,u5,u3} \\
                                        \text{服务器}\text{Xeon}
                                      \end{cases} \\
      \text{AMD(AMD)}
    \end{cases}\begin{cases}
      \text{核数/线程数:8,16,96}   \\
      \text{主频:2.5GHz,3.5GHz} \\
      \text{缓存(Cache)}
    \end{cases}
  \]
  Cerebras(CBRS),TBM
  \[
    \text{存储器}\begin{cases}
      \text{三星}        \\
      \text{海力士(SKHY)} \\
      \text{美光(MU)}
    \end{cases}\begin{cases}
      \text{容量:8GB,16GB}                  \\
      \text{频率:DDR4,DDR5,4000MHz,8000MHz} \\
      \text{时序:颗粒,A-die,B-die,C-die}
    \end{cases}
  \]
  套条(8GB*2,16GB*2) 优先插(2,4)槽
  \[
    \text{显卡}\begin{cases}
      \text{英伟达(NVDA)(MKL库)}\text{GTX}\rightarrow\text{RTX},\text{DLSS} \\
      \text{AMD}
    \end{cases}
  \]
\end{note}
\begin{note}
  软路由 openwrt AutoDL
\end{note}
\section{数学补充(DAY 3)}
\[
  \text{答案}\begin{cases}
    \text{判别问题(decision)} \\
    \text{优化问题(optimization)}
  \end{cases}
\]
\begin{exam}
  \[
    \text{判别问题}\begin{cases}
      \text{NP}\& \text{P} \\
      \text{NPC}           \\
      \text{NP-hard}
    \end{cases}
  \]
  N:non-deterministic（非确定性），P:Polynomial-time（多项式时间），C:Complete（完全）。
\end{exam}
\begin{defn}[P 类]
  能在多项式时间内被确定性算法判定的判别问题集合。
\end{defn}

\begin{defn}[NP 类]
  答案能被多项式时间内验证的问题集合；等价地，非确定性图灵机可在多项式时间内判定。显然 $P\subseteq NP$，$P=NP$ 是否成立是计算复杂性理论的核心未解问题。
\end{defn}

\begin{defn}[NP-hard]
  若 NP 中所有问题都能多项式时间归约到问题 $Q$，则 $Q$ 是 NP-hard——它至少和 NP 中任何问题一样难，但不要求自身属于 NP。
\end{defn}

\begin{defn}[NPC（NP 完全）]
  既属于 NP 又是 NP-hard 的问题。NPC 是 NP 中最"难"的一类：只要其中一个有多项式时间算法，则 $P=NP$。
\end{defn}
\begin{exam}
  \newterm{Karp 21 个 NPC 问题}：1972 年 Karp 证明了包括 3-SAT、团、顶点覆盖、哈密顿回路、背包等在内的 21 个问题都是 NP 完全的，奠定了 NP 完全性理论的基础。

  \newterm{3-SAT}：给定一个合取范式（CNF）公式，其中每个子句恰好含 3 个文字，问是否存在一组真值赋值使公式为真。3-SAT 是最经典的 NP 完全问题之一，大量 NPC 证明都从它出发做归约。

  \newterm{旅行商问题（TSP）}：见下节组合优化中的定义，是经典的 NP-hard 优化问题。
\end{exam}
\begin{defn}[哈密顿回路问题]
  给定一个无向图，判断是否存在一条经过每个顶点恰好一次的回路。
\end{defn}

\begin{defn}[欧拉路径问题]
  给定一个无向图，判断是否存在一条经过每条边恰好一次的路径。判定准则：奇数度的点个数为 0 或 2。当奇数度的点个数为 2 时，这 2 个奇数度的点为欧拉路径的起点和终点。
\end{defn}

\begin{defn}[旅行商问题]
  给定一个完全图，每条边都有一个权重（距离），要求找到一条最短的哈密顿回路。TSP 是 NP-hard 问题，没有已知的多项式时间算法可以解决所有实例。常用的近似算法包括贪心算法、动态规划、分支定界法和启发式算法（如遗传算法、模拟退火等）。
\end{defn}
\begin{exam}
  图 $\gG=(\gV,\gE)$，$\gV$为点(vertice)集，$\gE$为边(edge)集。上述三个问题是图论中最经典的路径/回路问题，其中哈密顿回路与 TSP 是 NP-hard，欧拉路径则存在多项式时间算法。

\end{exam}
\subsection{组合优化}
\[
  \text{优化问题}\begin{cases}
    \text{变量}   \\
    \text{目标函数} \\
    \text{约束条件}
  \end{cases}\begin{cases}
    \text{连续优化} \\
    \text{组合优化} \\
    \text{泛函优化} \\
  \end{cases}\begin{cases}
    \text{微分} \\
    \text{差分} \\
    \text{变分}
  \end{cases}
\]
\begin{exam}
  \[
    \begin{aligned}
      \min\quad        & f(x)                           \\
      \text{s.t.}\quad & f_i(x)\le 0,\quad i=1,\ldots,m \\
                       & h_j(x)=0,\quad j=1,\ldots,p
    \end{aligned}
  \]
\end{exam}
\begin{exam}
  $x$是一个排列
  \[
    \begin{aligned}
      \min\quad        & \sum_{i=1}^{n-1}w(x_i,x_{i+1})+w(x_n,x_1) \\
      \text{s.t.}\quad & (x_i,x_{i+1})\in \gE                      \\
                       & (x_n,x_1)\in\gE
    \end{aligned}
  \]
\end{exam}
\begin{defn}[近似算法（approximation algorithm）]
  在多项式时间内给出与最优解相差有界的可行解，用近似比衡量质量，例如 TSP 的 MST 2-近似算法、欧氏 TSP 的 Christofides 算法（近似比 $3/2$）。
\end{defn}

\begin{defn}[启发式算法（heuristic algorithm）]
  依靠经验规则快速找到可行解但不保证最优（如贪心），实现简单、速度快。
\end{defn}

\begin{defn}[仿生算法（bio-inspired / metaheuristic）]
  受自然过程启发的通用搜索框架，如遗传算法（GA）、粒子群优化（PSO）、蚁群算法（ACO）、模拟退火（SA）等，常用于大规模组合优化求近似最优。
\end{defn}

\begin{defn}[OPT 下界]
  能够多项式时间计算且不超过真实最优值 $\mathrm{OPT}$ 的界，用于评价近似算法的性能。例如最小生成树权值就是欧氏 TSP 的一个有效下界。
\end{defn}

\begin{defn}[最小生成树（MST）]
  在无向带权连通图中选出 $n-1$ 条边连通所有顶点且总权最小的树。Kruskal（贪心排序 + 并查集）伪代码如下：
  \begin{verbatim}
  function kruskal(V, edges):
      sort edges by weight ascending
      DSU = each vertex its own set
      T = []
      for (u, v, w) in sorted edges:
          if find(u) != find(v):
              union(u, v); T.append((u, v, w))
          if |T| == |V| - 1: break
      return T
  \end{verbatim}
  Prim 算法用堆每次取当前割边中权最小的边扩展生成树。两者复杂度均为 $O(m\log n)$。
\end{defn}

\begin{defn}[最大流问题]
  给定有向图、源点 $s$、汇点 $t$ 及边容量，求从 $s$ 到 $t$ 能输送的最大流量，等价于最小割（最大流最小割定理）；常用 Ford-Fulkerson、Edmonds-Karp、Dinic 算法求解。Ford-Fulkerson 伪代码如下：
  \begin{verbatim}
  function max_flow(G, s, t):
      f = 0
      while exists an augmenting path P in residual graph:
          c_f = min residual capacity on P
          augment f along P by c_f
          update residual capacities
      return f
  \end{verbatim}
\end{defn}
\begin{thm}[最大流最小割定理]
  在有向网络 $(V,E,c)$ 中，从 $s$ 到 $t$ 的最大流等于 $(s,t)$-割的最小容量。
\end{thm}
\begin{proof}
  一方面，任意可行流 $f$ 穿过任意 $(s,t)$-割 $(S,T)$ 时，由流量守恒与容量约束知 $|f| = f(S,T) \le c(S,T)$，故最大流 $\le$ 最小割。另一方面，Ford-Fulkerson 终止时残量网络中没有增广路，令 $S$ 为残量网络中从 $s$ 可达的顶点集，则 $s\in S$、$t\notin S$，且跨越 $(S,T)$ 的原边全部饱和、反向边全为 0，故 $|f|=c(S,T)$。两方向结合得等式。
\end{proof}
\begin{note}
  图 $\gG=(\gV,\gE)$，点集 $\gV=\{\gV_1,\ldots,\gV_n\}$，边集 $\gE=\{(u,v)\mid u,v\in\gV\}$，边权 $w(u,v)$。
\end{note}
\begin{defn}[度（degree）]
  与顶点关联的边数，记为 $\deg(v)$；无向图中所有顶点度之和等于边数的两倍。
\end{defn}

\begin{defn}[邻居（neighbor）]
  与 $u$ 有边直接相连的顶点集合 $N(u)=\{v\mid (u,v)\in\gE\}$。
\end{defn}

\begin{defn}[子图与生成子图]
  \newterm{子图}：顶点和边都是原图子集的图；\newterm{生成子图}：顶点集等于原图顶点集的子图。
\end{defn}

\begin{defn}[连通性与连通分量]
  \newterm{连通性}：无向图任意两点有路径则连通；有向图分为\textbf{强连通}（任意两点互相可达）、\textbf{弱连通}（忽略方向后连通）和\textbf{半连通}（任意两点至少一方可达）。\newterm{连通分量}（\textbf{极大连通子图}）：加上任意一个外部顶点都不再连通的连通子图，即无法再加入任何顶点仍保持连通的子图。
\end{defn}

\begin{defn}[树]
  连通且无环的无向图，$n$ 个点的树一定有 $n-1$ 条边，任意两点间有唯一路径；\newterm{最小生成树}见上文。
\end{defn}

为什么最小生成树可以成为 TSP 的一个很好的下界？TSP 最优回路去掉任意一条边后恰是一棵生成树，因此
\[
  \mathrm{OPT}_{\mathrm{TSP}} > \mathrm{OPT}_{\mathrm{MST}} > \mathrm{OPT}_{\min\{n-1\text{ 条边}\}},
\]
其中 $\mathrm{OPT}_{\mathrm{MST}}$ 可由多项式时间算法精确求出，从而给出 $\mathrm{OPT}_{\mathrm{TSP}}$ 的一个有效下界。
\section{模型与问题的差距(DAY4)}
\subsection{图论}
\begin{note}
  \[
    \text{树}\begin{cases}
      \text{链} \\
      \text{星}
    \end{cases}\rightarrow\text{直径}
  \]
\end{note}
\begin{exam}
  求树的最大直径？
  \[
    \begin{aligned}
      \max\quad        & \{\min d(u,v)\} \\
      \text{s.t.}\quad & u,v\in \gV
    \end{aligned}
  \]
  \newterm{思路}：树的直径定义为所有点对距离的最大值。经典做法是两次 DFS/BFS：任选一点 $s$ 出发找最远的点 $u$，再从 $u$ 出发找最远的点 $v$，则 $u,v$ 的距离就是直径。正确性基于"任意点到直径端点 $u$ 的距离是到整棵树最远的距离"这一性质，用反证法可证。

  \[
    \begin{aligned}
      \min\quad        & \left\{\max\limits_{v\neq u,v\in\gV}d(u,v)\right\} \\
      \text{s.t.}\quad & u\in\gV
    \end{aligned}
  \]
  这是求图中"最小偏心距"（图的半径对应的中心），可用多源最短路或 Floyd 求解。
\end{exam}
\begin{note}
  问题可行性与评价标准 benchmark?
\end{note}
\begin{defn}[最近公共祖先（LCA, Lowest Common Ancestor）]
  在一棵有根树中，两个节点 $u,v$ 的所有公共祖先里深度最大的那个。经典求法：
  \begin{itemize}
    \item \newterm{倍增法}：预处理每个节点的 $2^k$ 级祖先，先让 $u,v$ 对齐深度，再一起向上跳，$O(n\log n)$ 预处理、$O(\log n)$ 每次查询；
    \item \newterm{Tarjan 离线算法}：借助 DFS 和并查集一次遍历回答全部询问，$O(n+\alpha(n))$；
    \item 转化为欧拉序 + RMQ：用 ST 表支持 $O(1)$ 查询。
  \end{itemize}
  倍增法伪代码如下：
  \begin{verbatim}
  // 预处理: up[u][k] = u 的第 2^k 级祖先, depth[u] = u 的深度
  function lca(u, v):
      if depth[u] < depth[v]: swap(u, v)     // 对齐深度
      lift u up by (depth[u] - depth[v]) using up[][]
      if u == v: return u
      for k = log n down to 0:
          if up[u][k] != up[v][k]:
              u = up[u][k]; v = up[v][k]
      return up[u][0]
  \end{verbatim}
\end{defn}
\begin{note}
  一个无向无环图是一棵树。
\end{note}
\begin{note}
  图的路径：最短、最长、次长路。
\end{note}
\begin{defn}[匹配（matching）]
  两两不相邻（无公共顶点）的边集。
\end{defn}

\begin{defn}[最大匹配]
  边数最多的匹配；若匹配边覆盖所有顶点则称为\textbf{完美匹配}。
\end{defn}

\begin{defn}[匈牙利算法（Hungarian algorithm）]
  求解二分图最大匹配的经典算法，核心是不断寻找增广路并反转，直至找不到增广路；复杂度 $O(nm)$（$n$ 为点数，$m$ 为边数）。伪代码如下：
  \begin{verbatim}
  function hungarian(G):                     // G 为二分图
      match[v] = -1 for all right-side v
      for u in left side:
          visited = { }
          if not try_augment(u): break       // 找不到增广路
      return number of matched edges

  function try_augment(u):
      for each v adjacent to u:
          if v not visited:
              visited.insert(v)
              if match[v] == -1 or try_augment(match[v]):
                  match[v] = u; return True
      return False
  \end{verbatim}
\end{defn}

\begin{defn}[点覆盖（vertex cover）]
  顶点子集，使每条边至少有一个端点被选入；\newterm{最小覆盖}：顶点数最少的点覆盖。
\end{defn}

\begin{thm}[König 定理]
  在二分图中，最大匹配大小等于最小点覆盖大小。
\end{thm}
\begin{proof}
  记最大匹配为 $M$。先证 $|M|\le$ 任意点覆盖的大小：每条匹配边都必须被覆盖，而匹配边两两不相邻，故至少需 $|M|$ 个点，所以最小点覆盖 $\ge|M|$。再证存在大小为 $|M|$ 的点覆盖：从左侧未匹配点出发走"匹配边--非匹配边交替路"，记交替路上可达的左侧点集为 $L'$、可达的右侧点集为 $R$；令 $C=(L\setminus L')\cup R$。可验证 $C$ 是点覆盖且 $|C|=|M|$（这正是匈牙利算法终止时留下的结构），故最小点覆盖 $\le|M|$。  两方面合起来即得等式。
\end{proof}

参考：\href{https://oi-wiki.org}{oi-wiki.org}
\begin{lem}[Burnside 引理]
  设有限群 $G$ 作用在集合 $X$ 上，则 $X$ 在 $G$ 下的轨道数等于
  \[
    \frac{1}{|G|}\sum_{g\in G} |X^g|,
  \]
  其中 $X^g=\{x\in X\mid g(x)=x\}$ 为 $g$ 的不动点集合。
\end{lem}
\begin{proof}
  双计数集合 $A=\{(g,x)\mid g(x)=x\}$。按第一个分量计数得 $|A|=\sum_{g\in G}|X^g|$；按第二个分量计数，对每个轨道 $O$ 中的 $x$，满足 $g(x)=x$ 的 $g$ 恰构成一个陪集，个数均为 $|G|/|O|$，故每个轨道贡献 $|O|\cdot |G|/|O|=|G|$ 个对。因此 $|A|=(\text{轨道数})\cdot|G|$，整理即得结论。
\end{proof}

\begin{thm}[Pólya 计数定理]
  用 $k$ 种颜色给 $n$ 个位置染色，在置换群 $G$ 作用下本质不同的染色数为
  \[
    \frac{1}{|G|}\sum_{g\in G} k^{c(g)},
  \]
  其中 $c(g)$ 是置换 $g$ 中循环的个数。
\end{thm}
\begin{proof}
  由 Burnside 引理，轨道数 $=\frac1{|G|}\sum_{g\in G}|\text{Fix}(g)|$。一个染色在 $g$ 下不变的充要条件是其每个循环内的位置颜色相同：每个循环有 $k$ 种选色，且不同循环独立，故 $|\text{Fix}(g)|=k^{c(g)}$。代入即得。
\end{proof}

\begin{defn}[团（clique）]
  任意两点之间都有边相连的顶点子集（完全子图）；最大团问题是 NP-hard 的。
\end{defn}

\begin{defn}[染色问题]
  给图的每个顶点染色，使相邻顶点颜色不同，所用最少颜色数称为\textbf{色数}；判定图是否 3-可染色是 NP 完全问题。
\end{defn}
\begin{exam}
  有一个 $4\times 4$ 的矩阵，每个位置上都有一个灯泡。按下某个位置时，会改变它上下左右四个位置的开关状态。对任意给定的一种初始状态，问最少需要多少次操作才能使所有灯泡全部熄灭？

  \newterm{思路（熄灯问题 / Lights Out）}：设 $x_{ij}\in\{0,1\}$ 表示是否按下位置 $(i,j)$，则每个位置的最终状态由初始状态与它自身及上下左右共 $\le 5$ 个按钮的异或和决定。要求最终全灭，得到一个 $\mathbb{F}_2$ 上的线性方程组 $A\va=\vb$（$\vb$ 由初始状态决定），用高斯消元解出 $\va$。由于按两次等于没按，最少次数正是解中 $1$ 的个数；若解不唯一，取 $\sum x_{ij}$ 最小的那组。

  有一个 $m\times n$ 的矩阵，每个位置上可能存在一个怪物。可以选择某一行或某一列进行一次轰炸，问最少需要多少次才能消灭全部怪物？

  \newterm{思路（二分图最小点覆盖）}：把每一行与每一列分别作为二分图两侧的顶点，每个位于 $(i,j)$ 的怪物作为连接"第 $i$ 行"与"第 $j$ 列"的边。轰炸某行/列等价于在对应顶点中选点，目标是让每条边（怪物）至少被选中一次的一个端点覆盖——这正是\textbf{最小点覆盖}问题。由 König 定理，二分图中最小点覆盖大小等于最大匹配大小，故用匈牙利算法求最大匹配即可得到最少轰炸次数。
\end{exam}
\subsection{统计}
\begin{defn}[二项分布（Binomial distribution）]
  $n$ 次独立重复试验（每次成功概率为 $p$）中成功次数 $X$ 服从参数 $n,p$ 的二项分布 $X\sim B(n,p)$，其分布列为
  \[
    P(X=k)=\binom{n}{k}p^k(1-p)^{n-k},\qquad 0\le k\le n,
  \]
  期望 $\E[X]=np$，方差 $\Var(X)=np(1-p)$。
\end{defn}

\begin{defn}[泊松分布（Poisson distribution）]
  单位时间或单位区域内随机事件发生次数 $X$ 服从参数 $\lambda$ 的泊松分布 $X\sim \mathrm{Poi}(\lambda)$，其分布列为
  \[
    P(X=k)=\frac{\lambda^k e^{-\lambda}}{k!},\qquad k=0,1,2,\dots,
  \]
  期望与方差均为 $\lambda$。它是 $n$ 很大、$p$ 很小而 $np=\lambda$ 保持恒定时的二项分布极限（泊松近似）。
\end{defn}

\begin{defn}[指数分布（Exponential distribution）]
  连续型随机变量 $X$ 的密度函数为
  \[
    f(x)=\lambda e^{-\lambda x},\quad x\ge 0,
  \]
  记为 $X\sim \mathrm{Exp}(\lambda)$，期望 $\E[X]=1/\lambda$，方差 $\Var(X)=1/\lambda^2$，具有无记忆性 $P(X>s+t\mid X>s)=P(X>t)$。它是泊松过程两次事件之间的等待时间的分布。
\end{defn}
\end{document}
